\chapter{Open Reduction and Internal Fixation (ORIF)}
\section*{Overview}
ORIF surgically realigns a broken bone and secures it with plates, screws, rods, wires, or other implants.
\section*{Why This Test or Procedure Is Ordered}
It is used for displaced, unstable, open, intra-articular, or otherwise unsuitable-for-casting fractures.
\section*{How This Test or Procedure Is Performed}
The fracture is exposed or approached through limited incisions, reduced into alignment, and stabilized with selected hardware. Imaging confirms position.
\section*{Risks and Recovery}
Infection, bleeding, nerve or vessel injury, nonunion, malunion, stiffness, hardware failure, and post-traumatic arthritis are possible. Immobilization and rehabilitation depend on the bone and fixation.
\section*{References}
\begin{enumerate}\item MedlinePlus. Fracture Repair. U.S. National Library of Medicine; 2025.\item Current Medical Diagnosis and Treatment. McGraw-Hill; 2025.\end{enumerate}
\clearpage
